% page 422 du fac-simile (image p-421 du scan)
% bloc 200.7 x 168.5 mm ; marge gauche 3.0 mm
\pgc{-17.780mm}{0.000mm}{
\l{}
\l{}
\l{}
\l{}
\l{}
\l{}
\l{\cel{15}\sou{oxo}.\cel{2}(\sou{kemio)}.\cel{2}Korpo simpla, ek la aero, ed uzata por res-}
\l{\cel{18}piroe kombusto. -\cel{1}\sou{Simbolo kemial}a :\cel{2}O.}
\l{}
\l{\cel{15}\sou{oxalo}. (\sou{bot.)}\cel{3}Planto supala, ek la familio \textquotedbl{}poligonei\textquotedbl{},qua}
\l{\cel{18}saporas kelke acida, e quan onu uzas por facar supi o}
\l{\cel{18}kom legumo. - L.\cel{1}\sou{rumex acetosa}.- DEFIS.}
\l{}
\l{\cel{15}\sou{oxido.}\cel{2}Kompozuro de korpo simpla kun oxo, e qua ne su unio-}
\l{\cel{18}nas kun la bazi. - DEFIS.}
\l{}
\l{\cel{15}\sou{oxidazo}.\cel{2}(\sou{kemio biologiala}). Fermento dissolvebla en aquo}
\l{\cel{18}- kompozajo di qua un ek la kompozanti esas un molekulo}
\l{\cel{18}metala qua igas lu apta kaptar oxo e transportar ica adsur}
\l{\cel{18}substanco organika, oxidigebla facile. -}
\l{}
\l{\cel{15}\sou{oxigeno}. (\sou{kemio}). Oxo en la formo di gaso. - DEFIS.}
\l{}
\l{\cel{15}\sou{oxikoko}. (\sou{bot.}) Genero de erikacei - subarboreto reptera,}
\l{\cel{18}di qua la frukti - beri globetatra - redatra, acido-sapor-}
\l{\cel{18}oza, divenas kelke sukroza pos la komenco di la frosto-}
\l{\cel{18}periodo, ed uzata por facar konfituro. - L.\cel{1}\sou{veccinium}\cel{1}oxy-}
\l{\cel{18}\sou{coccos.}}
\l{}
\l{\cel{15}-\cel{1}\sou{oz-.}\cel{3}Sufixo qua implikas la ideo : \textquotedbl{}qua naturale havas...qua}
\l{\cel{18}kontenas naturale...\textquotedbl{} e \textquotedbl{}qua kontenas to quon indikas la ra-}
\l{\cel{18}diko verba\textquotedbl{} (sen impliko di anteeso, prezenteso o futureso;}
\l{\cel{18}ca nuanci indikesas per \textquotedbl{}-inta\textquotedbl{}, \textquotedbl{}-anta\textquotedbl{}, \textquotedbl{}-onta\textquotedbl{}.}
\l{}
\l{\cel{15}\sou{oziero}. (\sou{bot.}) Mikra speco di saliko, di qua la rami, flexebla,}
\l{\cel{18}uzesas por facar korbatraji. L.\cel{1}\sou{salix viminalis}. EF.}
\l{}
\l{\cel{15}\sou{ozono.}\cel{1}(kemio). Oxigeno kondensita.\cel{1}\sou{Simbolo kemiala}\cel{1}:}
\l{\cel{18}O3 (oxigeno) O2.}
}
